% ============================================================
%  REFERENCIAS BIBLIOGRÁFICAS
% ============================================================

\chapter*{REFERENCIAS BIBLIOGRÁFICAS}
\addcontentsline{toc}{chapter}{Referencias Bibliográficas}

% Completar con las referencias APA según avance la tesis.
% Ejemplo de entradas pendientes:

\begin{itemize}
    \item Alsentzer, E., Murphy, J. R., Boag, W., Weng, W.-H., Jin, D., Naumann, T., \& McDermott, M. (2019). Publicly available clinical BERT embeddings. \textit{arXiv preprint arXiv:1904.03323}.

    \item Chapman, W. W., Bridewell, W., Hanbury, P., Cooper, G. F., \& Buchanan, B. G. (2001). A simple algorithm for identifying negated findings and diseases in discharge summaries. \textit{Journal of Biomedical Informatics, 34}(5), 301--310.

    \item Lee, J., Yoon, W., Kim, S., Kim, D., Kim, S., So, C. H., \& Kang, J. (2020). BioBERT: a pre-trained biomedical language representation model for biomedical text mining. \textit{Bioinformatics, 36}(4), 1234--1240.

    \item Organización Mundial de la Salud. (2019). \textit{Acción mundial para la seguridad del paciente}. OMS.

    \item Ratner, A., De Sa, C., Wu, S., Selsam, D., \& Ré, C. (2017). Data programming: Creating large training sets, quickly. \textit{Advances in Neural Information Processing Systems, 29}.
\end{itemize}
